\QHeading{第1题　口算。}
\par\noindent


\begin{oralgrid}
\OralItem{1}{3.14+2.66}{5.8}
\OralItem{2}{0.75+7.16}{7.91}
\OralItem{3}{2.43+5.76}{8.19}
\OralItem{4}{1.08+1.29}{2.37}
\OralItem{5}{3.8+4.81}{8.61}
\OralItem{6}{1.54+4.36}{5.9}
\OralItem{7}{2.86+1.33}{4.19}
\OralItem{8}{2.65+6.25}{8.9}\end{oralgrid}
\par

\QHeading{口算思路补充（对应上表题号）}

\Answer{上面列出得数，下面按小题给出不同计算思路。计数单位法适用于小数，整数题可比较凑整、拆分或商不变等方法。}
\MethodTitle{第1小题：$3.14+2.66$}

\MethodBlock{口算思路一：按相同计数单位计算}{先看成若干个百分之一，按整数加减，最后还原成小数。}{\begin{aligned}[t]&\phantom{={}}3.14+2.66\\
&=(314+266)\div 100\\
&=580\div 100\\
&=5.8\end{aligned}}
\MethodTitle{第2小题：$0.75+7.16$}

\MethodBlock{口算思路一：按相同计数单位计算}{先看成若干个百分之一，按整数加减，最后还原成小数。}{\begin{aligned}[t]&\phantom{={}}0.75+7.16\\
&=(75+716)\div 100\\
&=791\div 100\\
&=7.91\end{aligned}}
\MethodTitle{第3小题：$2.43+5.76$}

\MethodBlock{口算思路一：按相同计数单位计算}{先看成若干个百分之一，按整数加减，最后还原成小数。}{\begin{aligned}[t]&\phantom{={}}2.43+5.76\\
&=(243+576)\div 100\\
&=819\div 100\\
&=8.19\end{aligned}}

\MethodBlock{口算思路二：凑整再减去多加的数}{把第二个数凑成整数，再减去多加的部分。}{\begin{aligned}[t]&\phantom{={}}2.43+5.76\\
&=2.43+6-0.24\\
&=8.43-0.24\\
&=8.19\end{aligned}}
\MethodTitle{第4小题：$1.08+1.29$}

\MethodBlock{口算思路一：按相同计数单位计算}{先看成若干个百分之一，按整数加减，最后还原成小数。}{\begin{aligned}[t]&\phantom{={}}1.08+1.29\\
&=(108+129)\div 100\\
&=237\div 100\\
&=2.37\end{aligned}}
\MethodTitle{第5小题：$3.8+4.81$}

\MethodBlock{口算思路一：按相同计数单位计算}{先看成若干个百分之一，按整数加减，最后还原成小数。}{\begin{aligned}[t]&\phantom{={}}3.8+4.81\\
&=(380+481)\div 100\\
&=861\div 100\\
&=8.61\end{aligned}}

\MethodBlock{口算思路二：凑整再减去多加的数}{把第二个数凑成整数，再减去多加的部分。}{\begin{aligned}[t]&\phantom{={}}3.8+4.81\\
&=3.8+5-0.19\\
&=8.8-0.19\\
&=8.61\end{aligned}}
\MethodTitle{第6小题：$1.54+4.36$}

\MethodBlock{口算思路一：按相同计数单位计算}{先看成若干个百分之一，按整数加减，最后还原成小数。}{\begin{aligned}[t]&\phantom{={}}1.54+4.36\\
&=(154+436)\div 100\\
&=590\div 100\\
&=5.9\end{aligned}}
\MethodTitle{第7小题：$2.86+1.33$}

\MethodBlock{口算思路一：按相同计数单位计算}{先看成若干个百分之一，按整数加减，最后还原成小数。}{\begin{aligned}[t]&\phantom{={}}2.86+1.33\\
&=(286+133)\div 100\\
&=419\div 100\\
&=4.19\end{aligned}}
\MethodTitle{第8小题：$2.65+6.25$}

\MethodBlock{口算思路一：按相同计数单位计算}{先看成若干个百分之一，按整数加减，最后还原成小数。}{\begin{aligned}[t]&\phantom{={}}2.65+6.25\\
&=(265+625)\div 100\\
&=890\div 100\\
&=8.9\end{aligned}}


\par\noindent

\QHeading{第2题　列竖式计算，带※的要验算。}

\par\noindent\makebox[\linewidth][c]{
\begin{tikzpicture}[x=1bp,y=-1bp,color=black,draw=black,text=black]
\begin{scope}\clip (76.0,223.0) rectangle (540.0,411.0);
\hdhead{86}{总复习}
\hdsection{1}{151}{口算。}
\hdsection{2}{299}{列竖式计算，带※的要验算。}
\hdsection{3}{699}{解决问题。}
\q{119}{197}{180}{12.5}{$3.14+2.66=$}
\q{119}{241}{180}{12.5}{$3.8+4.81=$}
\q{309}{197}{180}{12.5}{$0.75+7.16=$}
\q{309}{241}{180}{12.5}{$1.54+4.36=$}
\q{496}{197}{180}{12.5}{$2.43+5.76=$}
\q{496}{241}{180}{12.5}{$2.86+1.33=$}
\q{686}{197}{180}{12.5}{$1.08+1.29=$}
\q{686}{241}{180}{12.5}{$2.65+6.25=$}
\q{121}{344}{220}{12.5}{$5.84+7.29=$}
\q{374}{344}{220}{12.5}{$20.45-12.35=$}
\q{625}{344}{220}{12.5}{※$20.05-8.67=$}
\q{121}{515}{220}{12.5}{$9.3+15.76=$}
\q{374}{515}{220}{12.5}{$12.56-4.98=$}
\q{625}{515}{220}{12.5}{※$8.42+19.58=$}
\q{122}{743}{716}{12.5}{① 学校剪纸社团有同等规格的黄色纸和红色纸若干张，现在打算发给}
\q{122}{786}{716}{12.5}{　每位同学9张红纸，12张黄纸。发给24位同学以后，还剩112张。}
\q{122}{829}{716}{12.5}{　两种纸原来一共有多少张？}
\q{122}{956}{716}{12.5}{② 惠民超市准备购进大米、小米和高粱米三种粮食，要用载质量3吨的}
\q{122}{999}{716}{12.5}{　小货车一次运回惠民超市。已购买大米1.7吨，小米0.4吨，还可以}
\q{122}{1042}{716}{12.5}{　购买多少吨高粱米？}

\begin{scope}[color=answerred,draw=answerred,text=answerred]\ans{166.00}{128.67}{40.00}{10.5}{$5.8$}
\ans{166.00}{158.67}{40.00}{10.5}{$8.61$}
\ans{292.00}{128.67}{40.00}{10.5}{$7.91$}
\ans{292.00}{158.67}{40.00}{10.5}{$5.9$}
\ans{418.00}{128.67}{40.00}{10.5}{$8.19$}
\ans{418.00}{158.67}{40.00}{10.5}{$4.19$}
\ans{545.33}{128.67}{40.00}{10.5}{$2.37$}
\ans{545.33}{158.67}{40.00}{10.5}{$8.9$}
\ans{80.67}{254.00}{145.33}{11}{\begin{tikzpicture}[x=1em,y=1em,baseline=(current bounding box.center),draw=answerred,text=answerred]
\node[anchor=base,inner sep=0pt] at (0.800,0) {$5$};
\node[anchor=base,inner sep=0pt] at (1.600,0) {$.$};
\node[anchor=base,inner sep=0pt] at (1.900,0) {$8$};
\node[anchor=base,inner sep=0pt] at (2.700,0) {$4$};
\node[anchor=base,inner sep=0pt] at (0.800,-1.1) {$7$};
\node[anchor=base,inner sep=0pt] at (1.600,-1.1) {$.$};
\node[anchor=base,inner sep=0pt] at (1.900,-1.1) {$2$};
\node[anchor=base,inner sep=0pt] at (2.700,-1.1) {$9$};
\node[anchor=base,inner sep=0pt] at (0.000,-2.3) {$1$};
\node[anchor=base,inner sep=0pt] at (0.800,-2.3) {$3$};
\node[anchor=base,inner sep=0pt] at (1.600,-2.3) {$.$};
\node[anchor=base,inner sep=0pt] at (1.900,-2.3) {$1$};
\node[anchor=base,inner sep=0pt] at (2.700,-2.3) {$3$};
\node[anchor=base,inner sep=0pt] at (-.7,-1.1) {$+$};
\draw[line width=.5pt] (-1,-1.52)--(3.300,-1.52);
\node[anchor=base,inner sep=0pt,font=\normalsize] at (1.900,.72) {$1$};
\node[anchor=base,inner sep=0pt,font=\normalsize] at (0.800,.72) {$1$};
\node[anchor=base,inner sep=0pt,font=\normalsize] at (0.000,.72) {$1$};
\end{tikzpicture}}
\ans{249.33}{254.00}{145.33}{11}{\begin{tikzpicture}[x=1em,y=1em,baseline=(current bounding box.center),draw=answerred,text=answerred]
\node[anchor=base,inner sep=0pt] at (0.000,0) {$2$};
\node[anchor=base,inner sep=0pt] at (0.800,0) {$0$};
\node[anchor=base,inner sep=0pt] at (1.600,0) {$.$};
\node[anchor=base,inner sep=0pt] at (1.900,0) {$4$};
\node[anchor=base,inner sep=0pt] at (2.700,0) {$5$};
\node[anchor=base,inner sep=0pt] at (0.000,-1.1) {$1$};
\node[anchor=base,inner sep=0pt] at (0.800,-1.1) {$2$};
\node[anchor=base,inner sep=0pt] at (1.600,-1.1) {$.$};
\node[anchor=base,inner sep=0pt] at (1.900,-1.1) {$3$};
\node[anchor=base,inner sep=0pt] at (2.700,-1.1) {$5$};
\node[anchor=base,inner sep=0pt] at (0.800,-2.3) {$8$};
\node[anchor=base,inner sep=0pt] at (1.600,-2.3) {$.$};
\node[anchor=base,inner sep=0pt] at (1.900,-2.3) {$1$};
\node[anchor=base,inner sep=0pt] at (2.700,-2.3) {$0$};
\node[anchor=base,inner sep=0pt] at (-.7,-1.1) {$-$};
\draw[line width=.5pt] (-1,-1.52)--(3.300,-1.52);
\draw[line width=.35pt] (-0.160,.03)--(0.160,.43);
\node[anchor=base,inner sep=0pt,font=\normalsize] at (0.000,.72) {$1$};
\draw[line width=.35pt] (0.640,.03)--(0.960,.43);
\node[anchor=base,inner sep=0pt,font=\normalsize] at (0.800,.72) {$10$};
\end{tikzpicture}}
\ans{416.67}{254.00}{145.33}{9.8}{\begin{tikzpicture}[x=1em,y=1em,baseline=(current bounding box.center),draw=answerred,text=answerred]
\node[anchor=base,inner sep=0pt] at (0.000,0) {$2$};
\node[anchor=base,inner sep=0pt] at (0.800,0) {$0$};
\node[anchor=base,inner sep=0pt] at (1.600,0) {$.$};
\node[anchor=base,inner sep=0pt] at (1.900,0) {$0$};
\node[anchor=base,inner sep=0pt] at (2.700,0) {$5$};
\node[anchor=base,inner sep=0pt] at (0.800,-1.1) {$8$};
\node[anchor=base,inner sep=0pt] at (1.600,-1.1) {$.$};
\node[anchor=base,inner sep=0pt] at (1.900,-1.1) {$6$};
\node[anchor=base,inner sep=0pt] at (2.700,-1.1) {$7$};
\node[anchor=base,inner sep=0pt] at (0.000,-2.3) {$1$};
\node[anchor=base,inner sep=0pt] at (0.800,-2.3) {$1$};
\node[anchor=base,inner sep=0pt] at (1.600,-2.3) {$.$};
\node[anchor=base,inner sep=0pt] at (1.900,-2.3) {$3$};
\node[anchor=base,inner sep=0pt] at (2.700,-2.3) {$8$};
\node[anchor=base,inner sep=0pt] at (-.7,-1.1) {$-$};
\draw[line width=.5pt] (-1,-1.52)--(3.300,-1.52);
\draw[line width=.35pt] (-0.160,.03)--(0.160,.43);
\node[anchor=base,inner sep=0pt,font=\normalsize] at (0.000,.72) {$1$};
\draw[line width=.35pt] (0.640,.03)--(0.960,.43);
\node[anchor=base,inner sep=0pt,font=\normalsize] at (0.800,.72) {$9$};
\draw[line width=.35pt] (1.740,.03)--(2.060,.43);
\node[anchor=base,inner sep=0pt,font=\normalsize] at (1.900,.72) {$9$};
\draw[line width=.35pt] (2.540,.03)--(2.860,.43);
\node[anchor=base,inner sep=0pt,font=\normalsize] at (2.700,.72) {$15$};
\end{tikzpicture}\hspace{3pt}验算：\begin{tikzpicture}[x=1em,y=1em,baseline=(current bounding box.center),draw=answerred,text=answerred]
\node[anchor=base,inner sep=0pt] at (0.000,0) {$1$};
\node[anchor=base,inner sep=0pt] at (0.800,0) {$1$};
\node[anchor=base,inner sep=0pt] at (1.600,0) {$.$};
\node[anchor=base,inner sep=0pt] at (1.900,0) {$3$};
\node[anchor=base,inner sep=0pt] at (2.700,0) {$8$};
\node[anchor=base,inner sep=0pt] at (0.800,-1.1) {$8$};
\node[anchor=base,inner sep=0pt] at (1.600,-1.1) {$.$};
\node[anchor=base,inner sep=0pt] at (1.900,-1.1) {$6$};
\node[anchor=base,inner sep=0pt] at (2.700,-1.1) {$7$};
\node[anchor=base,inner sep=0pt] at (0.000,-2.3) {$2$};
\node[anchor=base,inner sep=0pt] at (0.800,-2.3) {$0$};
\node[anchor=base,inner sep=0pt] at (1.600,-2.3) {$.$};
\node[anchor=base,inner sep=0pt] at (1.900,-2.3) {$0$};
\node[anchor=base,inner sep=0pt] at (2.700,-2.3) {$5$};
\node[anchor=base,inner sep=0pt] at (-.7,-1.1) {$+$};
\draw[line width=.5pt] (-1,-1.52)--(3.300,-1.52);
\node[anchor=base,inner sep=0pt,font=\normalsize] at (1.900,.72) {$1$};
\node[anchor=base,inner sep=0pt,font=\normalsize] at (0.800,.72) {$1$};
\node[anchor=base,inner sep=0pt,font=\normalsize] at (0.000,.72) {$1$};
\end{tikzpicture}}
\ans{80.67}{368.00}{145.33}{11}{\begin{tikzpicture}[x=1em,y=1em,baseline=(current bounding box.center),draw=answerred,text=answerred]
\node[anchor=base,inner sep=0pt] at (0.800,0) {$9$};
\node[anchor=base,inner sep=0pt] at (1.600,0) {$.$};
\node[anchor=base,inner sep=0pt] at (1.900,0) {$3$};
\node[anchor=base,inner sep=0pt] at (2.700,0) {$0$};
\node[anchor=base,inner sep=0pt] at (0.000,-1.1) {$1$};
\node[anchor=base,inner sep=0pt] at (0.800,-1.1) {$5$};
\node[anchor=base,inner sep=0pt] at (1.600,-1.1) {$.$};
\node[anchor=base,inner sep=0pt] at (1.900,-1.1) {$7$};
\node[anchor=base,inner sep=0pt] at (2.700,-1.1) {$6$};
\node[anchor=base,inner sep=0pt] at (0.000,-2.3) {$2$};
\node[anchor=base,inner sep=0pt] at (0.800,-2.3) {$5$};
\node[anchor=base,inner sep=0pt] at (1.600,-2.3) {$.$};
\node[anchor=base,inner sep=0pt] at (1.900,-2.3) {$0$};
\node[anchor=base,inner sep=0pt] at (2.700,-2.3) {$6$};
\node[anchor=base,inner sep=0pt] at (-.7,-1.1) {$+$};
\draw[line width=.5pt] (-1,-1.52)--(3.300,-1.52);
\node[anchor=base,inner sep=0pt,font=\normalsize] at (0.800,.72) {$1$};
\node[anchor=base,inner sep=0pt,font=\normalsize] at (0.000,.72) {$1$};
\end{tikzpicture}}
\ans{249.33}{368.00}{145.33}{11}{\begin{tikzpicture}[x=1em,y=1em,baseline=(current bounding box.center),draw=answerred,text=answerred]
\node[anchor=base,inner sep=0pt] at (0.000,0) {$1$};
\node[anchor=base,inner sep=0pt] at (0.800,0) {$2$};
\node[anchor=base,inner sep=0pt] at (1.600,0) {$.$};
\node[anchor=base,inner sep=0pt] at (1.900,0) {$5$};
\node[anchor=base,inner sep=0pt] at (2.700,0) {$6$};
\node[anchor=base,inner sep=0pt] at (0.800,-1.1) {$4$};
\node[anchor=base,inner sep=0pt] at (1.600,-1.1) {$.$};
\node[anchor=base,inner sep=0pt] at (1.900,-1.1) {$9$};
\node[anchor=base,inner sep=0pt] at (2.700,-1.1) {$8$};
\node[anchor=base,inner sep=0pt] at (0.800,-2.3) {$7$};
\node[anchor=base,inner sep=0pt] at (1.600,-2.3) {$.$};
\node[anchor=base,inner sep=0pt] at (1.900,-2.3) {$5$};
\node[anchor=base,inner sep=0pt] at (2.700,-2.3) {$8$};
\node[anchor=base,inner sep=0pt] at (-.7,-1.1) {$-$};
\draw[line width=.5pt] (-1,-1.52)--(3.300,-1.52);
\draw[line width=.35pt] (-0.160,.03)--(0.160,.43);
\node[anchor=base,inner sep=0pt,font=\normalsize] at (0.000,.72) {$0$};
\draw[line width=.35pt] (0.640,.03)--(0.960,.43);
\node[anchor=base,inner sep=0pt,font=\normalsize] at (0.800,.72) {$11$};
\draw[line width=.35pt] (1.740,.03)--(2.060,.43);
\node[anchor=base,inner sep=0pt,font=\normalsize] at (1.900,.72) {$14$};
\draw[line width=.35pt] (2.540,.03)--(2.860,.43);
\node[anchor=base,inner sep=0pt,font=\normalsize] at (2.700,.72) {$16$};
\end{tikzpicture}}
\ans{416.67}{368.00}{145.33}{9.8}{\begin{tikzpicture}[x=1em,y=1em,baseline=(current bounding box.center),draw=answerred,text=answerred]
\node[anchor=base,inner sep=0pt] at (0.800,0) {$8$};
\node[anchor=base,inner sep=0pt] at (1.600,0) {$.$};
\node[anchor=base,inner sep=0pt] at (1.900,0) {$4$};
\node[anchor=base,inner sep=0pt] at (2.700,0) {$2$};
\node[anchor=base,inner sep=0pt] at (0.000,-1.1) {$1$};
\node[anchor=base,inner sep=0pt] at (0.800,-1.1) {$9$};
\node[anchor=base,inner sep=0pt] at (1.600,-1.1) {$.$};
\node[anchor=base,inner sep=0pt] at (1.900,-1.1) {$5$};
\node[anchor=base,inner sep=0pt] at (2.700,-1.1) {$8$};
\node[anchor=base,inner sep=0pt] at (0.000,-2.3) {$2$};
\node[anchor=base,inner sep=0pt] at (0.800,-2.3) {$8$};
\node[anchor=base,inner sep=0pt] at (1.600,-2.3) {$.$};
\node[anchor=base,inner sep=0pt] at (1.900,-2.3) {$0$};
\node[anchor=base,inner sep=0pt] at (2.700,-2.3) {$0$};
\node[anchor=base,inner sep=0pt] at (-.7,-1.1) {$+$};
\draw[line width=.5pt] (-1,-1.52)--(3.300,-1.52);
\node[anchor=base,inner sep=0pt,font=\normalsize] at (1.900,.72) {$1$};
\node[anchor=base,inner sep=0pt,font=\normalsize] at (0.800,.72) {$1$};
\node[anchor=base,inner sep=0pt,font=\normalsize] at (0.000,.72) {$1$};
\end{tikzpicture}\hspace{3pt}验算：\begin{tikzpicture}[x=1em,y=1em,baseline=(current bounding box.center),draw=answerred,text=answerred]
\node[anchor=base,inner sep=0pt] at (0.000,0) {$2$};
\node[anchor=base,inner sep=0pt] at (0.800,0) {$8$};
\node[anchor=base,inner sep=0pt] at (1.600,0) {$.$};
\node[anchor=base,inner sep=0pt] at (1.900,0) {$0$};
\node[anchor=base,inner sep=0pt] at (2.700,0) {$0$};
\node[anchor=base,inner sep=0pt] at (0.000,-1.1) {$1$};
\node[anchor=base,inner sep=0pt] at (0.800,-1.1) {$9$};
\node[anchor=base,inner sep=0pt] at (1.600,-1.1) {$.$};
\node[anchor=base,inner sep=0pt] at (1.900,-1.1) {$5$};
\node[anchor=base,inner sep=0pt] at (2.700,-1.1) {$8$};
\node[anchor=base,inner sep=0pt] at (0.800,-2.3) {$8$};
\node[anchor=base,inner sep=0pt] at (1.600,-2.3) {$.$};
\node[anchor=base,inner sep=0pt] at (1.900,-2.3) {$4$};
\node[anchor=base,inner sep=0pt] at (2.700,-2.3) {$2$};
\node[anchor=base,inner sep=0pt] at (-.7,-1.1) {$-$};
\draw[line width=.5pt] (-1,-1.52)--(3.300,-1.52);
\draw[line width=.35pt] (-0.160,.03)--(0.160,.43);
\node[anchor=base,inner sep=0pt,font=\normalsize] at (0.000,.72) {$1$};
\draw[line width=.35pt] (0.640,.03)--(0.960,.43);
\node[anchor=base,inner sep=0pt,font=\normalsize] at (0.800,.72) {$17$};
\draw[line width=.35pt] (1.740,.03)--(2.060,.43);
\node[anchor=base,inner sep=0pt,font=\normalsize] at (1.900,.72) {$9$};
\draw[line width=.35pt] (2.540,.03)--(2.860,.43);
\node[anchor=base,inner sep=0pt,font=\normalsize] at (2.700,.72) {$10$};
\end{tikzpicture}}
\ans{100.67}{570.00}{452.67}{10.5}{{\setlength{\jot}{-2pt}$\begin{aligned}[t]&\phantom{={}}(9+12)\times24+112\\&=21\times24+112\\&=504+112\\&=616\text{（张）}\end{aligned}$}}
\ans{100.67}{716.00}{452.67}{11}{{\setlength{\jot}{-2pt}$\begin{aligned}[t]&\phantom{={}}3-1.7-0.4\\&=3-(1.7+0.4)\\&=3-2.1\\&=0.9\text{（吨）}\end{aligned}$}\\答：还可以购买0.9吨高粱米。}
\ans{310.00}{612.00}{245.00}{10.5}{答：两种纸原来共有616张。}

\end{scope}
\end{scope}\pgfresetboundingbox\path[use as bounding box] (76.0,223.0) rectangle (540.0,411.0);
\end{tikzpicture}}\par



\par




\par\noindent

\Answer{
\QHeading{计算过程：2. 列竖式计算}

小数点对齐，缺位补零；从右往左逐列计算。下面完整列出进位、退位及各数位的计算过程。

\par\medskip\noindent
\textbf{第1小题：\(5.84+7.29=13.13\)}\par\smallskip
\noindent
\begin{tikzpicture}[x=1em,y=1em,baseline=(current bounding box.center),draw=answerred,text=answerred]
\node[anchor=base,inner sep=0pt] at (0.800,0) {$5$};
\node[anchor=base,inner sep=0pt] at (1.600,0) {$.$};
\node[anchor=base,inner sep=0pt] at (1.900,0) {$8$};
\node[anchor=base,inner sep=0pt] at (2.700,0) {$4$};
\node[anchor=base,inner sep=0pt] at (0.800,-1.1) {$7$};
\node[anchor=base,inner sep=0pt] at (1.600,-1.1) {$.$};
\node[anchor=base,inner sep=0pt] at (1.900,-1.1) {$2$};
\node[anchor=base,inner sep=0pt] at (2.700,-1.1) {$9$};
\node[anchor=base,inner sep=0pt] at (0.000,-2.3) {$1$};
\node[anchor=base,inner sep=0pt] at (0.800,-2.3) {$3$};
\node[anchor=base,inner sep=0pt] at (1.600,-2.3) {$.$};
\node[anchor=base,inner sep=0pt] at (1.900,-2.3) {$1$};
\node[anchor=base,inner sep=0pt] at (2.700,-2.3) {$3$};
\node[anchor=base,inner sep=0pt] at (-.7,-1.1) {$+$};
\draw[line width=.5pt] (-1,-1.52)--(3.300,-1.52);
\node[anchor=base,inner sep=0pt,font=\normalsize] at (1.900,.72) {$1$};
\node[anchor=base,inner sep=0pt,font=\normalsize] at (0.800,.72) {$1$};
\node[anchor=base,inner sep=0pt,font=\normalsize] at (0.000,.72) {$1$};
\end{tikzpicture}
\par\smallskip{\normalsize 上方小字表示从后一位进来的数。\par}
\vspace{0pt}
\begin{enumerate}[leftmargin=1.7em,itemsep=5pt,topsep=0pt]
\item 百分位：\(\begin{aligned}[t]&\phantom{={}}4+9\\&=13\end{aligned}\)\quad 本位写3，向前一位进1。
\item 十分位：\(\begin{aligned}[t]&\phantom{={}}8+2+1\\&=10+1\\&=11\end{aligned}\)\quad 本位写1，向前一位进1。
\item 个位：\(\begin{aligned}[t]&\phantom{={}}5+7+1\\&=12+1\\&=13\end{aligned}\)\quad 本位写3，向前一位进1。
\end{enumerate}
最高位再写进上来的1。\par
\par\smallskip
因此，\(\begin{aligned}[t]&\phantom{={}}5.84+7.29\\&=13.13\end{aligned}\)。
\par\medskip

\par\medskip\noindent
\textbf{第2小题：\(20.45-12.35=8.1\)}\par\smallskip
\noindent
\begin{tikzpicture}[x=1em,y=1em,baseline=(current bounding box.center),draw=answerred,text=answerred]
\node[anchor=base,inner sep=0pt] at (0.000,0) {$2$};
\node[anchor=base,inner sep=0pt] at (0.800,0) {$0$};
\node[anchor=base,inner sep=0pt] at (1.600,0) {$.$};
\node[anchor=base,inner sep=0pt] at (1.900,0) {$4$};
\node[anchor=base,inner sep=0pt] at (2.700,0) {$5$};
\node[anchor=base,inner sep=0pt] at (0.000,-1.1) {$1$};
\node[anchor=base,inner sep=0pt] at (0.800,-1.1) {$2$};
\node[anchor=base,inner sep=0pt] at (1.600,-1.1) {$.$};
\node[anchor=base,inner sep=0pt] at (1.900,-1.1) {$3$};
\node[anchor=base,inner sep=0pt] at (2.700,-1.1) {$5$};
\node[anchor=base,inner sep=0pt] at (0.800,-2.3) {$8$};
\node[anchor=base,inner sep=0pt] at (1.600,-2.3) {$.$};
\node[anchor=base,inner sep=0pt] at (1.900,-2.3) {$1$};
\node[anchor=base,inner sep=0pt] at (2.700,-2.3) {$0$};
\node[anchor=base,inner sep=0pt] at (-.7,-1.1) {$-$};
\draw[line width=.5pt] (-1,-1.52)--(3.300,-1.52);
\draw[line width=.35pt] (-0.160,.03)--(0.160,.43);
\node[anchor=base,inner sep=0pt,font=\normalsize] at (0.000,.72) {$1$};
\draw[line width=.35pt] (0.640,.03)--(0.960,.43);
\node[anchor=base,inner sep=0pt,font=\normalsize] at (0.800,.72) {$10$};
\end{tikzpicture}
\par\smallskip{\normalsize 划线数字是退位前的数，上方小字是退位后参与计算的数。\par}
\vspace{0pt}
\begin{enumerate}[leftmargin=1.7em,itemsep=5pt,topsep=0pt]
\item 百分位：\(\begin{aligned}[t]&\phantom{={}}5-5\\&=0\end{aligned}\)\quad 本位写0。
\item 十分位：\(\begin{aligned}[t]&\phantom{={}}4-3\\&=1\end{aligned}\)\quad 本位写1。
\item 十位退位：\(\begin{aligned}[t]&\phantom{={}}2-1\\&=1\end{aligned}\)\quad 借出1给后一位。
\item 个位接收退位：\(\begin{aligned}[t]&\phantom{={}}0+10\\&=10\end{aligned}\)\quad 前一位的1变成本位的10。
\item 个位：\(\begin{aligned}[t]&\phantom{={}}10-2\\&=8\end{aligned}\)\quad 本位写8。
\item 十位：\(\begin{aligned}[t]&\phantom{={}}1-1\\&=0\end{aligned}\)\quad 本位写0。
\end{enumerate}
\par\smallskip
因此，\(\begin{aligned}[t]&\phantom{={}}20.45-12.35\\&=8.10\end{aligned}\)。
\par\medskip

\par\medskip\noindent
\textbf{第3小题：\(20.05-8.67=11.38\)}\par\smallskip
\noindent
\begin{tikzpicture}[x=1em,y=1em,baseline=(current bounding box.center),draw=answerred,text=answerred]
\node[anchor=base,inner sep=0pt] at (0.000,0) {$2$};
\node[anchor=base,inner sep=0pt] at (0.800,0) {$0$};
\node[anchor=base,inner sep=0pt] at (1.600,0) {$.$};
\node[anchor=base,inner sep=0pt] at (1.900,0) {$0$};
\node[anchor=base,inner sep=0pt] at (2.700,0) {$5$};
\node[anchor=base,inner sep=0pt] at (0.800,-1.1) {$8$};
\node[anchor=base,inner sep=0pt] at (1.600,-1.1) {$.$};
\node[anchor=base,inner sep=0pt] at (1.900,-1.1) {$6$};
\node[anchor=base,inner sep=0pt] at (2.700,-1.1) {$7$};
\node[anchor=base,inner sep=0pt] at (0.000,-2.3) {$1$};
\node[anchor=base,inner sep=0pt] at (0.800,-2.3) {$1$};
\node[anchor=base,inner sep=0pt] at (1.600,-2.3) {$.$};
\node[anchor=base,inner sep=0pt] at (1.900,-2.3) {$3$};
\node[anchor=base,inner sep=0pt] at (2.700,-2.3) {$8$};
\node[anchor=base,inner sep=0pt] at (-.7,-1.1) {$-$};
\draw[line width=.5pt] (-1,-1.52)--(3.300,-1.52);
\draw[line width=.35pt] (-0.160,.03)--(0.160,.43);
\node[anchor=base,inner sep=0pt,font=\normalsize] at (0.000,.72) {$1$};
\draw[line width=.35pt] (0.640,.03)--(0.960,.43);
\node[anchor=base,inner sep=0pt,font=\normalsize] at (0.800,.72) {$9$};
\draw[line width=.35pt] (1.740,.03)--(2.060,.43);
\node[anchor=base,inner sep=0pt,font=\normalsize] at (1.900,.72) {$9$};
\draw[line width=.35pt] (2.540,.03)--(2.860,.43);
\node[anchor=base,inner sep=0pt,font=\normalsize] at (2.700,.72) {$15$};
\end{tikzpicture}
\par\smallskip{\normalsize 划线数字是退位前的数，上方小字是退位后参与计算的数。\par}
\vspace{0pt}
\begin{enumerate}[leftmargin=1.7em,itemsep=5pt,topsep=0pt]
\item 十位退位：\(\begin{aligned}[t]&\phantom{={}}2-1\\&=1\end{aligned}\)\quad 借出1给后一位。
\item 个位传递退位：\(\begin{aligned}[t]&\phantom{={}}0+10-1\\&=10-1\\&=9\end{aligned}\)\quad 收到10，再借出1给后一位。
\item 十分位传递退位：\(\begin{aligned}[t]&\phantom{={}}0+10-1\\&=10-1\\&=9\end{aligned}\)\quad 收到10，再借出1给后一位。
\item 百分位接收退位：\(\begin{aligned}[t]&\phantom{={}}5+10\\&=15\end{aligned}\)\quad 前一位的1变成本位的10。
\item 百分位：\(\begin{aligned}[t]&\phantom{={}}15-7\\&=8\end{aligned}\)\quad 本位写8。
\item 十分位：\(\begin{aligned}[t]&\phantom{={}}9-6\\&=3\end{aligned}\)\quad 本位写3。
\item 个位：\(\begin{aligned}[t]&\phantom{={}}9-8\\&=1\end{aligned}\)\quad 本位写1。
\item 十位：\(\begin{aligned}[t]&\phantom{={}}1-0\\&=1\end{aligned}\)\quad 本位写1。
\end{enumerate}
\par\smallskip
因此，\(\begin{aligned}[t]&\phantom{={}}20.05-8.67\\&=11.38\end{aligned}\)。
\par\medskip

\par\medskip\noindent
\textbf{第3小题验算：\(11.38+8.67\)}\par\smallskip
\noindent
\begin{tikzpicture}[x=1em,y=1em,baseline=(current bounding box.center),draw=answerred,text=answerred]
\node[anchor=base,inner sep=0pt] at (0.000,0) {$1$};
\node[anchor=base,inner sep=0pt] at (0.800,0) {$1$};
\node[anchor=base,inner sep=0pt] at (1.600,0) {$.$};
\node[anchor=base,inner sep=0pt] at (1.900,0) {$3$};
\node[anchor=base,inner sep=0pt] at (2.700,0) {$8$};
\node[anchor=base,inner sep=0pt] at (0.800,-1.1) {$8$};
\node[anchor=base,inner sep=0pt] at (1.600,-1.1) {$.$};
\node[anchor=base,inner sep=0pt] at (1.900,-1.1) {$6$};
\node[anchor=base,inner sep=0pt] at (2.700,-1.1) {$7$};
\node[anchor=base,inner sep=0pt] at (0.000,-2.3) {$2$};
\node[anchor=base,inner sep=0pt] at (0.800,-2.3) {$0$};
\node[anchor=base,inner sep=0pt] at (1.600,-2.3) {$.$};
\node[anchor=base,inner sep=0pt] at (1.900,-2.3) {$0$};
\node[anchor=base,inner sep=0pt] at (2.700,-2.3) {$5$};
\node[anchor=base,inner sep=0pt] at (-.7,-1.1) {$+$};
\draw[line width=.5pt] (-1,-1.52)--(3.300,-1.52);
\node[anchor=base,inner sep=0pt,font=\normalsize] at (1.900,.72) {$1$};
\node[anchor=base,inner sep=0pt,font=\normalsize] at (0.800,.72) {$1$};
\node[anchor=base,inner sep=0pt,font=\normalsize] at (0.000,.72) {$1$};
\end{tikzpicture}
\par\smallskip{\normalsize 上方小字表示从后一位进来的数。\par}
\vspace{0pt}
\begin{enumerate}[leftmargin=1.7em,itemsep=5pt,topsep=0pt]
\item 百分位：\(\begin{aligned}[t]&\phantom{={}}8+7\\&=15\end{aligned}\)\quad 本位写5，向前一位进1。
\item 十分位：\(\begin{aligned}[t]&\phantom{={}}3+6+1\\&=9+1\\&=10\end{aligned}\)\quad 本位写0，向前一位进1。
\item 个位：\(\begin{aligned}[t]&\phantom{={}}1+8+1\\&=9+1\\&=10\end{aligned}\)\quad 本位写0，向前一位进1。
\item 十位：\(\begin{aligned}[t]&\phantom{={}}1+0+1\\&=1+1\\&=2\end{aligned}\)\quad 本位写2。
\end{enumerate}
\par\smallskip
因此，\(\begin{aligned}[t]&\phantom{={}}11.38+8.67\\&=20.05\end{aligned}\)。
\par\medskip

\par\medskip\noindent
\textbf{第4小题：\(9.3+15.76=25.06\)}\par\smallskip
\noindent
\begin{tikzpicture}[x=1em,y=1em,baseline=(current bounding box.center),draw=answerred,text=answerred]
\node[anchor=base,inner sep=0pt] at (0.800,0) {$9$};
\node[anchor=base,inner sep=0pt] at (1.600,0) {$.$};
\node[anchor=base,inner sep=0pt] at (1.900,0) {$3$};
\node[anchor=base,inner sep=0pt] at (2.700,0) {$0$};
\node[anchor=base,inner sep=0pt] at (0.000,-1.1) {$1$};
\node[anchor=base,inner sep=0pt] at (0.800,-1.1) {$5$};
\node[anchor=base,inner sep=0pt] at (1.600,-1.1) {$.$};
\node[anchor=base,inner sep=0pt] at (1.900,-1.1) {$7$};
\node[anchor=base,inner sep=0pt] at (2.700,-1.1) {$6$};
\node[anchor=base,inner sep=0pt] at (0.000,-2.3) {$2$};
\node[anchor=base,inner sep=0pt] at (0.800,-2.3) {$5$};
\node[anchor=base,inner sep=0pt] at (1.600,-2.3) {$.$};
\node[anchor=base,inner sep=0pt] at (1.900,-2.3) {$0$};
\node[anchor=base,inner sep=0pt] at (2.700,-2.3) {$6$};
\node[anchor=base,inner sep=0pt] at (-.7,-1.1) {$+$};
\draw[line width=.5pt] (-1,-1.52)--(3.300,-1.52);
\node[anchor=base,inner sep=0pt,font=\normalsize] at (0.800,.72) {$1$};
\node[anchor=base,inner sep=0pt,font=\normalsize] at (0.000,.72) {$1$};
\end{tikzpicture}
\par\smallskip{\normalsize 上方小字表示从后一位进来的数。\par}
\vspace{0pt}
\begin{enumerate}[leftmargin=1.7em,itemsep=5pt,topsep=0pt]
\item 百分位：\(\begin{aligned}[t]&\phantom{={}}0+6\\&=6\end{aligned}\)\quad 本位写6。
\item 十分位：\(\begin{aligned}[t]&\phantom{={}}3+7\\&=10\end{aligned}\)\quad 本位写0，向前一位进1。
\item 个位：\(\begin{aligned}[t]&\phantom{={}}9+5+1\\&=14+1\\&=15\end{aligned}\)\quad 本位写5，向前一位进1。
\item 十位：\(\begin{aligned}[t]&\phantom{={}}0+1+1\\&=1+1\\&=2\end{aligned}\)\quad 本位写2。
\end{enumerate}
\par\smallskip
因此，\(\begin{aligned}[t]&\phantom{={}}9.3+15.76\\&=25.06\end{aligned}\)。
\par\medskip

\par\medskip\noindent
\textbf{第5小题：\(12.56-4.98=7.58\)}\par\smallskip
\noindent
\begin{tikzpicture}[x=1em,y=1em,baseline=(current bounding box.center),draw=answerred,text=answerred]
\node[anchor=base,inner sep=0pt] at (0.000,0) {$1$};
\node[anchor=base,inner sep=0pt] at (0.800,0) {$2$};
\node[anchor=base,inner sep=0pt] at (1.600,0) {$.$};
\node[anchor=base,inner sep=0pt] at (1.900,0) {$5$};
\node[anchor=base,inner sep=0pt] at (2.700,0) {$6$};
\node[anchor=base,inner sep=0pt] at (0.800,-1.1) {$4$};
\node[anchor=base,inner sep=0pt] at (1.600,-1.1) {$.$};
\node[anchor=base,inner sep=0pt] at (1.900,-1.1) {$9$};
\node[anchor=base,inner sep=0pt] at (2.700,-1.1) {$8$};
\node[anchor=base,inner sep=0pt] at (0.800,-2.3) {$7$};
\node[anchor=base,inner sep=0pt] at (1.600,-2.3) {$.$};
\node[anchor=base,inner sep=0pt] at (1.900,-2.3) {$5$};
\node[anchor=base,inner sep=0pt] at (2.700,-2.3) {$8$};
\node[anchor=base,inner sep=0pt] at (-.7,-1.1) {$-$};
\draw[line width=.5pt] (-1,-1.52)--(3.300,-1.52);
\draw[line width=.35pt] (-0.160,.03)--(0.160,.43);
\node[anchor=base,inner sep=0pt,font=\normalsize] at (0.000,.72) {$0$};
\draw[line width=.35pt] (0.640,.03)--(0.960,.43);
\node[anchor=base,inner sep=0pt,font=\normalsize] at (0.800,.72) {$11$};
\draw[line width=.35pt] (1.740,.03)--(2.060,.43);
\node[anchor=base,inner sep=0pt,font=\normalsize] at (1.900,.72) {$14$};
\draw[line width=.35pt] (2.540,.03)--(2.860,.43);
\node[anchor=base,inner sep=0pt,font=\normalsize] at (2.700,.72) {$16$};
\end{tikzpicture}
\par\smallskip{\normalsize 划线数字是退位前的数，上方小字是退位后参与计算的数。\par}
\vspace{0pt}
\begin{enumerate}[leftmargin=1.7em,itemsep=5pt,topsep=0pt]
\item 十分位退位：\(\begin{aligned}[t]&\phantom{={}}5-1\\&=4\end{aligned}\)\quad 借出1给后一位。
\item 百分位接收退位：\(\begin{aligned}[t]&\phantom{={}}6+10\\&=16\end{aligned}\)\quad 前一位的1变成本位的10。
\item 百分位：\(\begin{aligned}[t]&\phantom{={}}16-8\\&=8\end{aligned}\)\quad 本位写8。
\item 个位退位：\(\begin{aligned}[t]&\phantom{={}}2-1\\&=1\end{aligned}\)\quad 借出1给后一位。
\item 十分位接收退位：\(\begin{aligned}[t]&\phantom{={}}4+10\\&=14\end{aligned}\)\quad 前一位的1变成本位的10。
\item 十分位：\(\begin{aligned}[t]&\phantom{={}}14-9\\&=5\end{aligned}\)\quad 本位写5。
\item 十位退位：\(\begin{aligned}[t]&\phantom{={}}1-1\\&=0\end{aligned}\)\quad 借出1给后一位。
\item 个位接收退位：\(\begin{aligned}[t]&\phantom{={}}1+10\\&=11\end{aligned}\)\quad 前一位的1变成本位的10。
\item 个位：\(\begin{aligned}[t]&\phantom{={}}11-4\\&=7\end{aligned}\)\quad 本位写7。
\item 十位：\(\begin{aligned}[t]&\phantom{={}}0-0\\&=0\end{aligned}\)\quad 本位写0。
\end{enumerate}
\par\smallskip
因此，\(\begin{aligned}[t]&\phantom{={}}12.56-4.98\\&=7.58\end{aligned}\)。
\par\medskip

\par\medskip\noindent
\textbf{第6小题：\(8.42+19.58=28\)}\par\smallskip
\noindent
\begin{tikzpicture}[x=1em,y=1em,baseline=(current bounding box.center),draw=answerred,text=answerred]
\node[anchor=base,inner sep=0pt] at (0.800,0) {$8$};
\node[anchor=base,inner sep=0pt] at (1.600,0) {$.$};
\node[anchor=base,inner sep=0pt] at (1.900,0) {$4$};
\node[anchor=base,inner sep=0pt] at (2.700,0) {$2$};
\node[anchor=base,inner sep=0pt] at (0.000,-1.1) {$1$};
\node[anchor=base,inner sep=0pt] at (0.800,-1.1) {$9$};
\node[anchor=base,inner sep=0pt] at (1.600,-1.1) {$.$};
\node[anchor=base,inner sep=0pt] at (1.900,-1.1) {$5$};
\node[anchor=base,inner sep=0pt] at (2.700,-1.1) {$8$};
\node[anchor=base,inner sep=0pt] at (0.000,-2.3) {$2$};
\node[anchor=base,inner sep=0pt] at (0.800,-2.3) {$8$};
\node[anchor=base,inner sep=0pt] at (1.600,-2.3) {$.$};
\node[anchor=base,inner sep=0pt] at (1.900,-2.3) {$0$};
\node[anchor=base,inner sep=0pt] at (2.700,-2.3) {$0$};
\node[anchor=base,inner sep=0pt] at (-.7,-1.1) {$+$};
\draw[line width=.5pt] (-1,-1.52)--(3.300,-1.52);
\node[anchor=base,inner sep=0pt,font=\normalsize] at (1.900,.72) {$1$};
\node[anchor=base,inner sep=0pt,font=\normalsize] at (0.800,.72) {$1$};
\node[anchor=base,inner sep=0pt,font=\normalsize] at (0.000,.72) {$1$};
\end{tikzpicture}
\par\smallskip{\normalsize 上方小字表示从后一位进来的数。\par}
\vspace{0pt}
\begin{enumerate}[leftmargin=1.7em,itemsep=5pt,topsep=0pt]
\item 百分位：\(\begin{aligned}[t]&\phantom{={}}2+8\\&=10\end{aligned}\)\quad 本位写0，向前一位进1。
\item 十分位：\(\begin{aligned}[t]&\phantom{={}}4+5+1\\&=9+1\\&=10\end{aligned}\)\quad 本位写0，向前一位进1。
\item 个位：\(\begin{aligned}[t]&\phantom{={}}8+9+1\\&=17+1\\&=18\end{aligned}\)\quad 本位写8，向前一位进1。
\item 十位：\(\begin{aligned}[t]&\phantom{={}}0+1+1\\&=1+1\\&=2\end{aligned}\)\quad 本位写2。
\end{enumerate}
\par\smallskip
因此，\(\begin{aligned}[t]&\phantom{={}}8.42+19.58\\&=28.00\end{aligned}\)。
\par\medskip

\par\medskip\noindent
\textbf{第6小题验算：\(28.00-19.58\)}\par\smallskip
\noindent
\begin{tikzpicture}[x=1em,y=1em,baseline=(current bounding box.center),draw=answerred,text=answerred]
\node[anchor=base,inner sep=0pt] at (0.000,0) {$2$};
\node[anchor=base,inner sep=0pt] at (0.800,0) {$8$};
\node[anchor=base,inner sep=0pt] at (1.600,0) {$.$};
\node[anchor=base,inner sep=0pt] at (1.900,0) {$0$};
\node[anchor=base,inner sep=0pt] at (2.700,0) {$0$};
\node[anchor=base,inner sep=0pt] at (0.000,-1.1) {$1$};
\node[anchor=base,inner sep=0pt] at (0.800,-1.1) {$9$};
\node[anchor=base,inner sep=0pt] at (1.600,-1.1) {$.$};
\node[anchor=base,inner sep=0pt] at (1.900,-1.1) {$5$};
\node[anchor=base,inner sep=0pt] at (2.700,-1.1) {$8$};
\node[anchor=base,inner sep=0pt] at (0.800,-2.3) {$8$};
\node[anchor=base,inner sep=0pt] at (1.600,-2.3) {$.$};
\node[anchor=base,inner sep=0pt] at (1.900,-2.3) {$4$};
\node[anchor=base,inner sep=0pt] at (2.700,-2.3) {$2$};
\node[anchor=base,inner sep=0pt] at (-.7,-1.1) {$-$};
\draw[line width=.5pt] (-1,-1.52)--(3.300,-1.52);
\draw[line width=.35pt] (-0.160,.03)--(0.160,.43);
\node[anchor=base,inner sep=0pt,font=\normalsize] at (0.000,.72) {$1$};
\draw[line width=.35pt] (0.640,.03)--(0.960,.43);
\node[anchor=base,inner sep=0pt,font=\normalsize] at (0.800,.72) {$17$};
\draw[line width=.35pt] (1.740,.03)--(2.060,.43);
\node[anchor=base,inner sep=0pt,font=\normalsize] at (1.900,.72) {$9$};
\draw[line width=.35pt] (2.540,.03)--(2.860,.43);
\node[anchor=base,inner sep=0pt,font=\normalsize] at (2.700,.72) {$10$};
\end{tikzpicture}
\par\smallskip{\normalsize 划线数字是退位前的数，上方小字是退位后参与计算的数。\par}
\vspace{0pt}
\begin{enumerate}[leftmargin=1.7em,itemsep=5pt,topsep=0pt]
\item 个位退位：\(\begin{aligned}[t]&\phantom{={}}8-1\\&=7\end{aligned}\)\quad 借出1给后一位。
\item 十分位传递退位：\(\begin{aligned}[t]&\phantom{={}}0+10-1\\&=10-1\\&=9\end{aligned}\)\quad 收到10，再借出1给后一位。
\item 百分位接收退位：\(\begin{aligned}[t]&\phantom{={}}0+10\\&=10\end{aligned}\)\quad 前一位的1变成本位的10。
\item 百分位：\(\begin{aligned}[t]&\phantom{={}}10-8\\&=2\end{aligned}\)\quad 本位写2。
\item 十分位：\(\begin{aligned}[t]&\phantom{={}}9-5\\&=4\end{aligned}\)\quad 本位写4。
\item 十位退位：\(\begin{aligned}[t]&\phantom{={}}2-1\\&=1\end{aligned}\)\quad 借出1给后一位。
\item 个位接收退位：\(\begin{aligned}[t]&\phantom{={}}7+10\\&=17\end{aligned}\)\quad 前一位的1变成本位的10。
\item 个位：\(\begin{aligned}[t]&\phantom{={}}17-9\\&=8\end{aligned}\)\quad 本位写8。
\item 十位：\(\begin{aligned}[t]&\phantom{={}}1-1\\&=0\end{aligned}\)\quad 本位写0。
\end{enumerate}
\par\smallskip
因此，\(\begin{aligned}[t]&\phantom{={}}28.00-19.58\\&=8.42\end{aligned}\)。
\par\medskip

}

\QHeading{第3题　解决问题。}

\Problem{① 学校剪纸社团有同等规格的黄色纸和红色纸若干张，现在打算发给每位同学9张红纸，12张黄纸。发给24位同学以后，还剩112张。两种纸原来一共有多少张？}
\par


\MethodTitle{方法一}
\Answer{{\Formula{\begin{aligned}[t]&\phantom{={}}(9+12)\times24+112\\&=21\times24+112\\&=504+112\\&=616\text{（张）}\end{aligned}}}}
\Answer{答：两种纸原来共有616张。}

\MethodTitle{方法二：分别求发出的红纸、黄纸，再加余下的}
\Answer{先分别计算24位同学各领到的红纸与黄纸总量。}
\Answer{\Formula{\begin{aligned}[t]&\phantom{={}}9\times 24+12\times 24+112\\
&=216+288+112\\
&=504+112\\
&=616\text{（张）}\end{aligned}}}
\Answer{答：两种纸原来一共有616张。}
\Problem{② 惠民超市准备购进大米、小米和高粱米三种粮食，要用载质量3吨的小货车一次运回惠民超市。已购买大米1.7吨，小米0.4吨，还可以购买多少吨高粱米？}

\MethodTitle{方法一}
\Answer{{\Formula{\begin{aligned}[t]&\phantom{={}}3-1.7-0.4\\&=3-(1.7+0.4)\\&=3-2.1\\&=0.9\text{（吨）}\end{aligned}}}\par 答：还可以购买0.9吨高粱米。}

\MethodTitle{方法二：按装车顺序连续减}
\Answer{先扣去大米的质量，再扣去小米的质量。}
\Answer{\Formula{\begin{aligned}[t]&\phantom{={}}3-1.7-0.4\\
&=1.3-0.4\\
&=0.9\text{（吨）}\end{aligned}}}
\Answer{答：还可以购买0.9吨高粱米。}

